% !TEX TS-program = lualatex
% !TEX encoding = UTF-8 Unicode
\documentclass[12pt,a4paper]{article}

\IfFileExists{src/libs/body.tex}{% --- body.tex ---
\usepackage{pdfpages}
\usepackage{fontspec}
% Prefer Times New Roman when available; fall back in container environments.
\IfFontExistsTF{Times New Roman}{
  \setmainfont{Times New Roman}
}{
  \setmainfont{Liberation Serif}
}
\usepackage[vietnamese]{babel}

% Page layout and spacing
\usepackage[left=3cm,right=2cm,top=2cm,bottom=2cm]{geometry}
\usepackage{titlesec}
\usepackage{setspace}

% Silence noisy warnings
\usepackage{silence}
\WarningFilter{transparent}{Loading aborted}
\WarningFilter{hyperref}{The anchor of a bookmark and its parent's}

% Core packages
\usepackage{float}
\usepackage{svg}
\usepackage{graphicx}
\graphicspath{{assets/}{../assets/}{../../assets/}}
\svgpath{{assets/}{../assets/}{../../assets/}}
\usepackage{enumitem}
\usepackage{array}
\usepackage{tabularx}
\usepackage{multirow}
\usepackage[table]{xcolor}
\usepackage{ulem}
\usepackage{xurl}

% ====== THÊM 2 GÓI NÀY ĐỂ FIX LỖI ADJUSTBOX VÀ FOREST ======
\usepackage{adjustbox}
\usepackage[edges]{forest}
% ============================================================

% TikZ
\usepackage{tikz}

% ====== BỔ SUNG FIT, BACKGROUNDS, TREES VÀO DANH SÁCH ======
\usetikzlibrary{calc, arrows.meta, shapes.geometric, positioning, fit, backgrounds, trees}

% Table tuning
\renewcommand{\tabularxcolumn}[1]{m{#1}}
\hbadness=10000
\hfuzz=10000pt

% Shared commands
\newcommand{\subsecheading}[1]{%
    \par\vspace{2pt}%
    \hspace{1.5em}\textbf{#1}%
    \par\vspace{2pt}%
}
\newcommand{\introtext}[1]{%
    \par\hspace{3.0em}#1\par%
}
\newcommand{\StandaloneDocumentSetup}[2]{%
  \pagenumbering{arabic}%
  \ApplyStandardFooter{#1}{#2}%
}
\newcommand{\StandaloneSectionSetup}[2]{%
  \ApplyStandardFooter{#1}{#2}%
}

% Math
\usepackage{amsmath}

% Hyperref
\usepackage{hyperref}
\hypersetup{
    colorlinks=true,
    linkcolor=black,
    filecolor=magenta,
    urlcolor=blue,
}}{%
  \IfFileExists{../libs/body.tex}{% --- body.tex ---
\usepackage{pdfpages}
\usepackage{fontspec}
% Prefer Times New Roman when available; fall back in container environments.
\IfFontExistsTF{Times New Roman}{
  \setmainfont{Times New Roman}
}{
  \setmainfont{Liberation Serif}
}
\usepackage[vietnamese]{babel}

% Page layout and spacing
\usepackage[left=3cm,right=2cm,top=2cm,bottom=2cm]{geometry}
\usepackage{titlesec}
\usepackage{setspace}

% Silence noisy warnings
\usepackage{silence}
\WarningFilter{transparent}{Loading aborted}
\WarningFilter{hyperref}{The anchor of a bookmark and its parent's}

% Core packages
\usepackage{float}
\usepackage{svg}
\usepackage{graphicx}
\graphicspath{{assets/}{../assets/}{../../assets/}}
\svgpath{{assets/}{../assets/}{../../assets/}}
\usepackage{enumitem}
\usepackage{array}
\usepackage{tabularx}
\usepackage{multirow}
\usepackage[table]{xcolor}
\usepackage{ulem}
\usepackage{xurl}

% ====== THÊM 2 GÓI NÀY ĐỂ FIX LỖI ADJUSTBOX VÀ FOREST ======
\usepackage{adjustbox}
\usepackage[edges]{forest}
% ============================================================

% TikZ
\usepackage{tikz}

% ====== BỔ SUNG FIT, BACKGROUNDS, TREES VÀO DANH SÁCH ======
\usetikzlibrary{calc, arrows.meta, shapes.geometric, positioning, fit, backgrounds, trees}

% Table tuning
\renewcommand{\tabularxcolumn}[1]{m{#1}}
\hbadness=10000
\hfuzz=10000pt

% Shared commands
\newcommand{\subsecheading}[1]{%
    \par\vspace{2pt}%
    \hspace{1.5em}\textbf{#1}%
    \par\vspace{2pt}%
}
\newcommand{\introtext}[1]{%
    \par\hspace{3.0em}#1\par%
}
\newcommand{\StandaloneDocumentSetup}[2]{%
  \pagenumbering{arabic}%
  \ApplyStandardFooter{#1}{#2}%
}
\newcommand{\StandaloneSectionSetup}[2]{%
  \ApplyStandardFooter{#1}{#2}%
}

% Math
\usepackage{amsmath}

% Hyperref
\usepackage{hyperref}
\hypersetup{
    colorlinks=true,
    linkcolor=black,
    filecolor=magenta,
    urlcolor=blue,
}}{% --- body.tex ---
\usepackage{pdfpages}
\usepackage{fontspec}
% Prefer Times New Roman when available; fall back in container environments.
\IfFontExistsTF{Times New Roman}{
  \setmainfont{Times New Roman}
}{
  \setmainfont{Liberation Serif}
}
\usepackage[vietnamese]{babel}

% Page layout and spacing
\usepackage[left=3cm,right=2cm,top=2cm,bottom=2cm]{geometry}
\usepackage{titlesec}
\usepackage{setspace}

% Silence noisy warnings
\usepackage{silence}
\WarningFilter{transparent}{Loading aborted}
\WarningFilter{hyperref}{The anchor of a bookmark and its parent's}

% Core packages
\usepackage{float}
\usepackage{svg}
\usepackage{graphicx}
\graphicspath{{assets/}{../assets/}{../../assets/}}
\svgpath{{assets/}{../assets/}{../../assets/}}
\usepackage{enumitem}
\usepackage{array}
\usepackage{tabularx}
\usepackage{multirow}
\usepackage[table]{xcolor}
\usepackage{ulem}
\usepackage{xurl}

% ====== THÊM 2 GÓI NÀY ĐỂ FIX LỖI ADJUSTBOX VÀ FOREST ======
\usepackage{adjustbox}
\usepackage[edges]{forest}
% ============================================================

% TikZ
\usepackage{tikz}

% ====== BỔ SUNG FIT, BACKGROUNDS, TREES VÀO DANH SÁCH ======
\usetikzlibrary{calc, arrows.meta, shapes.geometric, positioning, fit, backgrounds, trees}

% Table tuning
\renewcommand{\tabularxcolumn}[1]{m{#1}}
\hbadness=10000
\hfuzz=10000pt

% Shared commands
\newcommand{\subsecheading}[1]{%
    \par\vspace{2pt}%
    \hspace{1.5em}\textbf{#1}%
    \par\vspace{2pt}%
}
\newcommand{\introtext}[1]{%
    \par\hspace{3.0em}#1\par%
}
\newcommand{\StandaloneDocumentSetup}[2]{%
  \pagenumbering{arabic}%
  \ApplyStandardFooter{#1}{#2}%
}
\newcommand{\StandaloneSectionSetup}[2]{%
  \ApplyStandardFooter{#1}{#2}%
}

% Math
\usepackage{amsmath}

% Hyperref
\usepackage{hyperref}
\hypersetup{
    colorlinks=true,
    linkcolor=black,
    filecolor=magenta,
    urlcolor=blue,
}}%
}
\IfFileExists{src/libs/header.tex}{% --- header.tex ---
\usepackage{fancyhdr}
\setlength{\headheight}{15.5pt} % Thêm dòng này để fix cảnh báo chiều cao
%\setlength{\headheight}{14.5pt}
\addtolength{\topmargin}{-2.5pt}

\pagestyle{fancy}
\fancyhf{}
\renewcommand{\headrulewidth}{0pt}
\renewcommand{\footrulewidth}{0pt}

\newcommand{\ApplyStandardHeader}{%
  \fancyhead[L]{%
    \begin{minipage}[t]{0.795\linewidth}%
      UIT \textbar{} SS004.F21 \textbar{} Nhóm 03 \textbar{} Đồ án cuối kỳ%
    \end{minipage}%
    \vrule%
    \begin{minipage}[t]{0.185\linewidth}%
      \centering cTetris%
    \end{minipage}%
  }%
}

\ApplyStandardHeader
}{%
  \IfFileExists{../libs/header.tex}{% --- header.tex ---
\usepackage{fancyhdr}
\setlength{\headheight}{15.5pt} % Thêm dòng này để fix cảnh báo chiều cao
%\setlength{\headheight}{14.5pt}
\addtolength{\topmargin}{-2.5pt}

\pagestyle{fancy}
\fancyhf{}
\renewcommand{\headrulewidth}{0pt}
\renewcommand{\footrulewidth}{0pt}

\newcommand{\ApplyStandardHeader}{%
  \fancyhead[L]{%
    \begin{minipage}[t]{0.795\linewidth}%
      UIT \textbar{} SS004.F21 \textbar{} Nhóm 03 \textbar{} Đồ án cuối kỳ%
    \end{minipage}%
    \vrule%
    \begin{minipage}[t]{0.185\linewidth}%
      \centering cTetris%
    \end{minipage}%
  }%
}

\ApplyStandardHeader
}{% --- header.tex ---
\usepackage{fancyhdr}
\setlength{\headheight}{15.5pt} % Thêm dòng này để fix cảnh báo chiều cao
%\setlength{\headheight}{14.5pt}
\addtolength{\topmargin}{-2.5pt}

\pagestyle{fancy}
\fancyhf{}
\renewcommand{\headrulewidth}{0pt}
\renewcommand{\footrulewidth}{0pt}

\newcommand{\ApplyStandardHeader}{%
  \fancyhead[L]{%
    \begin{minipage}[t]{0.795\linewidth}%
      UIT \textbar{} SS004.F21 \textbar{} Nhóm 03 \textbar{} Đồ án cuối kỳ%
    \end{minipage}%
    \vrule%
    \begin{minipage}[t]{0.185\linewidth}%
      \centering cTetris%
    \end{minipage}%
  }%
}

\ApplyStandardHeader
}%
}
\IfFileExists{src/libs/footer.tex}{% --- footer.tex ---
\usepackage{lastpage}

\newcommand{\ApplyStandardFooter}[2]{%
  \fancyfoot[L]{%
    \begin{minipage}[t]{0.795\linewidth}%
      Document ID: #1%
    \end{minipage}%
    \vrule%
    \begin{minipage}[t]{0.185\linewidth}%
      \centering\thepage/\pageref{#2}%
    \end{minipage}%
  }%
  \fancypagestyle{plain}{%
    \fancyhf{}%
    \renewcommand{\headrulewidth}{0pt}%
    \renewcommand{\footrulewidth}{0pt}%
    \ApplyStandardHeader%
    \fancyfoot[L]{%
      \begin{minipage}[t]{0.795\linewidth}%
        Document ID: #1%
      \end{minipage}%
      \vrule%
      \begin{minipage}[t]{0.185\linewidth}%
        \centering\thepage/\pageref{#2}%
      \end{minipage}%
    }%
  }%
  \pagestyle{fancy}%
}
}{%
  \IfFileExists{../libs/footer.tex}{% --- footer.tex ---
\usepackage{lastpage}

\newcommand{\ApplyStandardFooter}[2]{%
  \fancyfoot[L]{%
    \begin{minipage}[t]{0.795\linewidth}%
      Document ID: #1%
    \end{minipage}%
    \vrule%
    \begin{minipage}[t]{0.185\linewidth}%
      \centering\thepage/\pageref{#2}%
    \end{minipage}%
  }%
  \fancypagestyle{plain}{%
    \fancyhf{}%
    \renewcommand{\headrulewidth}{0pt}%
    \renewcommand{\footrulewidth}{0pt}%
    \ApplyStandardHeader%
    \fancyfoot[L]{%
      \begin{minipage}[t]{0.795\linewidth}%
        Document ID: #1%
      \end{minipage}%
      \vrule%
      \begin{minipage}[t]{0.185\linewidth}%
        \centering\thepage/\pageref{#2}%
      \end{minipage}%
    }%
  }%
  \pagestyle{fancy}%
}
}{% --- footer.tex ---
\usepackage{lastpage}

\newcommand{\ApplyStandardFooter}[2]{%
  \fancyfoot[L]{%
    \begin{minipage}[t]{0.795\linewidth}%
      Document ID: #1%
    \end{minipage}%
    \vrule%
    \begin{minipage}[t]{0.185\linewidth}%
      \centering\thepage/\pageref{#2}%
    \end{minipage}%
  }%
  \fancypagestyle{plain}{%
    \fancyhf{}%
    \renewcommand{\headrulewidth}{0pt}%
    \renewcommand{\footrulewidth}{0pt}%
    \ApplyStandardHeader%
    \fancyfoot[L]{%
      \begin{minipage}[t]{0.795\linewidth}%
        Document ID: #1%
      \end{minipage}%
      \vrule%
      \begin{minipage}[t]{0.185\linewidth}%
        \centering\thepage/\pageref{#2}%
      \end{minipage}%
    }%
  }%
  \pagestyle{fancy}%
}
}%
}
\usepackage{docmute}


\hypersetup{pdftitle={Đặc tả Thiết kế Hệ thống - Công nghệ nền tảng}}



\begin{document}
\providecommand{\StandaloneSectionSetup}[2]{\ApplyStandardFooter{#1}{#2}}
\StandaloneSectionSetup{07-gameTechnologies.tex}{LastPageTechnologyAdoption}

\onehalfspacing

\section{Công nghệ Nền tảng}

% ── GIAO DIỆN NGƯỜI DÙNG ────────────────────────────────────────────────────
\subsection{Công nghệ Giao diện Người dùng (Frontend)}
\subsubsection{Tổng quan Kiến trúc}
Hệ thống tuân theo nguyên tắc "Một mã nguồn, đa nền tảng": Mã C++ chuẩn 17 được biên dịch cho cả hệ điều hành máy tính (macOS/Windows/Ubuntu) và Trình duyệt Web (WebAssembly).

\begin{center}
\begin{adjustbox}{max width=\linewidth}
\begin{tikzpicture}[node distance=0.4cm]
  \node[techBox=tBlue, text width=12cm]
    (entry) at (0,0)
    {Thư viện SDL3 --- Cửa sổ · Bộ vẽ đồ họa · Vòng lặp sự kiện · Luồng âm thanh};

  \node[comp] (svg)  at (-4.5,-1.8) {Bộ chuyển đổi\\Ảnh Vector\\Thành điểm ảnh};
  \node[comp] (draw) at (-1.5,-1.8) {Vẽ đồ họa\\bằng thuật toán\\(Không ảnh tĩnh)};
  \node[comp] (dbg)  at (1.5,-1.8)  {Vẽ chữ\\(Phông chữ\\tích hợp)};
  \node[comp] (aud)  at (4.5,-1.8)  {Luồng âm thanh\\\& Điều chỉnh\\âm lượng};

  \node[techBox=tOrange, text width=12cm]
    (wasm) at (0,-3.4)
    {Nền tảng WebAssembly --- Trình biên dịch · Hàm dị bộ · Tệp ảo (IDBFS) · Trang nền};

  \foreach \c in {svg,draw,dbg,aud}{
    \draw[arrow](entry)--(\c.north);
    \draw[arrow](\c.south)--(wasm.north -| \c);
  }
\end{tikzpicture}
\end{adjustbox}
\end{center}

\subsubsection{Thư viện Giao diện}
\begin{center}
\begin{adjustbox}{max width=\linewidth}
\begin{tabular}{|p{3.5cm}|p{2cm}|p{9.5cm}|}
  \hline
  \rowcolor{gray!20}
  \textbf{Thư viện} & \textbf{Phiên bản} & \textbf{Vai trò cốt lõi} \\
  \hline
  \textbf{SDL3} & Mới nhất & Cửa sổ, bộ vẽ 2D, vòng lặp sự kiện, luồng âm thanh. Hỗ trợ đa nền tảng. \\
  \hline
  \textbf{SDL3\_ttf} & Tích hợp sẵn & Hiển thị chữ cho hội thoại; dùng hàm vẽ chữ gốc nếu thiếu phông. \\
  \hline
  \textbf{nanosvg(rast)} & Tệp duy nhất & Phân tích và chuyển đổi hình ảnh vector sang kết cấu đồ họa (khởi tạo lười 1 lần). \\
  \hline
  \textbf{Emscripten} & 3.1.x & Biên dịch C++ thành WebAssembly. Cho phép vòng lặp đồng bộ qua cờ \texttt{ASYNCIFY}. \\
  \hline
  \textbf{Hệ thống Tệp ảo} & Trình duyệt & Ánh xạ đường dẫn lưu trữ thành cơ sở dữ liệu trình duyệt IndexedDB. \\
  \hline
  \textbf{Trang nền HTML} & Tùy chỉnh & Khung chứa Web, đặt ở giữa, gắn biểu tượng và tiêu đề bản quyền. \\
  \hline
\end{tabular}
\end{adjustbox}
\end{center}

\subsubsection{Quy trình xuất Khung hình}
\begin{center}
\begin{adjustbox}{max width=\linewidth}
\begin{tikzpicture}[
  b/.style={techBoxSm=tBlue, text width=2.2cm},
  a/.style={sdsArrow=tBlue}
]
  \node[b](p1){Bắt Sự kiện\\(Lớp đầu vào)};
  \node[b,right=0.5cm of p1](p2){Cập nhật\\Trạng thái};
  \node[b,right=0.5cm of p2](p3){Xóa\\Màn hình};
  \node[b,right=0.5cm of p3](p4){Vẽ Hình nền\\(Chất liệu SVG)};
  \node[b,right=0.5cm of p4](p5){Vẽ Phân hệ\\(Bảng/Giao diện)};
  \node[b,right=0.5cm of p5](p6){Xuất hình\\(60 FPS)};
  \foreach \x/\y in {p1/p2,p2/p3,p3/p4,p4/p5,p5/p6} {\draw[a](\x)--(\y);}
\end{tikzpicture}
\end{adjustbox}
\end{center}
\textit{Tỉ lệ màn hình cố định 9:16 (270$\times$480 px), không phản hồi (responsive) nhằm tránh sai lệch bố cục giữa Máy tính và Web.}

% ── MÁY CHỦ & DỮ LIỆU (BACKEND) ─────────────────────────────────────────────
\subsection{Công nghệ Máy chủ \& Dữ liệu (Backend)}

\subsubsection{Kiến trúc Tổng thể}
"Máy chủ" bao gồm Động cơ C++ chạy cục bộ kết hợp kênh phân phối (Máy chủ Gist) và xếp hạng toàn cầu (Cloudflare D1).
\begin{center}
\begin{adjustbox}{max width=\linewidth}
\begin{tikzpicture}[
  box/.style={techBoxSm=tGreen, text width=2.6cm},
  cloud/.style={techBoxSm=tOrange, text width=2.6cm},
  arrow/.style={sdsArrow=tGreen},
  darrow/.style={sdsArrow=tOrange,dashed}
]
  \node[techBox=tGreen, text width=12cm]
    (core) at (0,0)
    {Động cơ trò chơi C++17 --- Cốt truyện · Bảng điều khiển · Lõi trò chơi};

  \node[box](gs)  at (-4.5,-1.8) {Đồng bộ Truyện\\\& Vòng lặp\\Hội thoại};
  \node[box](gc)  at (-1.5,-1.8) {Trung tâm ĐK\\+ Cài đặt\\+ Chọn Truyện};
  \node[box](gr)  at (1.5,-1.8)  {Vòng lặp Vật lý\\+ Ghi Lưu trữ};
  \node[box](http) at (4.5,-1.8) {Lớp Mạng HTTP\\(Máy tính \& Web)};

  \node[cloud](gist)   at (-4.5,-4.8) {Dịch vụ Gist\\(Khai báo \&\\Nội dung truyện)};
  \node[cloud](d1)     at (0.0,-4.8)  {Đám mây D1\\(Cơ sở dữ liệu\\Xếp hạng)};
  \node[cloud](cfwork) at (4.5,-4.8)  {Máy chủ Trung gian\\(Xử lý Yêu cầu)};

  \foreach \c in {gs,gc,gr,http}{\draw[arrow](core)--(\c.north);}
  
  \draw[darrow](gs.south)--(gist.north) node[midway,right,font=\tiny]{Mã băm SHA};
  \draw[darrow](http.south)--(cfwork.north) node[midway,right,font=\tiny]{Gửi thành tích};
  \draw[darrow](cfwork.west)--(d1.east) node[midway,above,font=\tiny]{Ghi dữ liệu};

  \node[font=\tiny,vTag=orange] at (-4.5,-5.6){Bản 2};
  \node[font=\tiny,vTag=red]   at (0.0,-5.6){Bản 3};
  \node[font=\tiny,vTag=red]   at (4.5,-5.6){Bản 3};
\end{tikzpicture}
\end{adjustbox}
\end{center}

\subsubsection{Chi tiết Công nghệ Máy chủ}
\begin{center}
\begin{adjustbox}{max width=\linewidth}
\begin{tabular}{|p{3cm}|p{2cm}|p{10cm}|}
  \hline
  \rowcolor{gray!20}
  \textbf{Công nghệ} & \textbf{Phiên bản} & \textbf{Vai trò} \\
  \hline
  \textbf{C++ Chuẩn 17} & Đa nền tảng & Ngôn ngữ chính; không cấp phát động trong vòng lặp game; dọn dẹp kết nối tự động. \\
  \hline
  \textbf{Phân tích JSON} & Tệp tiêu đề & Đọc dữ liệu khai báo và nội dung chương truyện từ JSON; bảo vệ lỗi ngoại lệ. \\
  \hline
  \textbf{Mạng HTTP} & Hệ thống \& Web & Thư viện \texttt{libcurl} cho máy tính, hàm tải gốc cho Web (chặn đồng bộ khi khởi động). \\
  \hline
  \textbf{Máy chủ Gist} & API v3 & Kênh phân phối nội dung: so sánh Mã băm để chỉ tải dữ liệu mới. \\
  \hline
  \textbf{Máy chủ Cloudflare} & Bản 3 & Máy chủ trung gian nhận điểm số thành tích và lưu vào CSDL. \\
  \hline
  \textbf{Đám mây D1} & Bản 3 & CSDL xếp hạng toàn cầu. Cục bộ là bản chính, đám mây là bản sao xem chung. \\
  \hline
\end{tabular}
\end{adjustbox}
\end{center}

\subsubsection{Luồng đồng bộ nội dung (Mã băm SHA)}
\begin{center}
\begin{adjustbox}{max width=\linewidth}
\begin{tikzpicture}[
  proc/.style={techBoxSm=tGreen, text width=3.0cm, minimum height=0.7cm},
  dec/.style={diamond, draw=tOrange, fill=tOrange!8, aspect=2.4, align=center, font=\scriptsize, inner sep=1pt},
  arr/.style={sdsArrow=tGreen}
]
  \node[proc](p0){Tải Khai báo\\từ Máy chủ};
  \node[proc,right=0.6cm of p0](p1){Đọc Phiên bản\\Mới nhất};
  \node[dec,right=0.6cm of p1](d0){Mã băm\\Khớp?};
  \node[proc,right=0.6cm of d0](p2){Bỏ qua\\Chương};
  \node[proc,below=0.8cm of d0](p3){Tải Dữ liệu\\Chương (JSON)};
  \node[proc,left=0.6cm of p3](p4){\texttt{ApDungDuLieu()}\\+ CẬP NHẬT DB};
  \node[proc,left=0.6cm of p4](p5){\texttt{LuuMaBam()}\\+ Đồng bộ Tệp Ảo};

  \draw[arr](p0)--(p1); \draw[arr](p1)--(d0);
  \draw[arr](d0)--node[above,font=\tiny]{Có}(p2);
  \draw[arr](d0)--node[right,font=\tiny]{Không}(p3);
  \draw[arr](p3)--(p4); \draw[arr](p4)--(p5);
  \draw[arr,gray,dashed](p2.south) |- ([yshift=-0.5cm]p5.south) -- (p5.south);
\end{tikzpicture}
\end{adjustbox}
\end{center}

% ── LƯU TRỮ (DATABASE) ──────────────────────────────────────────────────────
\subsection{Kiến trúc Lưu trữ}
Hệ thống kết hợp SQLite3 nhúng (Cục bộ), Tệp ảo (Trình duyệt), và Đám mây D1 (Toàn cầu).

\begin{center}
\begin{adjustbox}{max width=\linewidth}
\begin{tikzpicture}[
  grp/.style={rectangle, draw=#1!60, fill=#1!8, rounded corners=4pt, thick, align=center, font=\scriptsize, inner sep=6pt},
  tbl/.style={rectangle, draw=#1!50, fill=#1!5, rounded corners=2pt, align=left, font=\scriptsize, inner sep=4pt, text width=3.0cm},
  arrow/.style={sdsArrow=gray!60, dashed}
]
  \node[grp=blue, minimum width=3.6cm, minimum height=8.0cm] (g1box) at (-4.5,0) {};
  \node[font=\scriptsize\bfseries\color{blue}, align=center] at ([yshift=-0.85cm]g1box.south) {Nhóm 1 — Đọc/Ghi\\Người dùng};
  
  \node[tbl=blue] (rec) at (-4.5, 2.5) {\textbf{Bảng\_LichSu}\\ID, NguoiDung\\TGBatDau, TGKetThuc\\DiemSo, ThoiGian\\TocDo, LanChoi};
  \node[tbl=blue] (sto) at (-4.5, 0.1) {\textbf{Bảng\_TienDo}\\NguoiDung, Truyen\\KichHoat, DaChon\\DiemCao, TocDoCao};
  \node[tbl=blue] (set) at (-4.5, -2.4) {\textbf{Bảng\_CaiDat}\\Khoa\_ID (Chính)\\GiaTri};

  \node[grp=purple, minimum width=3.6cm, minimum height=8.0cm] (g2box) at (0,0) {};
  \node[font=\scriptsize\bfseries\color{purple}, align=center] at (0, -4.85) {Nhóm 2 — Dữ liệu Dùng chung};

  \node[tbl=purple] (sd) at (0, 2.8) {\textbf{Bảng\_DuLieu}\\Truyen, Chuong\\Ten, NguongDiem\\NguongTocDo, MaTran};
  \node[tbl=purple] (sdlg) at (0, 1.0) {\textbf{Bảng\_HoiThoai}\\Truyen, Chuong, Nut\\NhanVat, LoiThoai\\Anh, NutKeTiep};
  \node[tbl=purple] (scho) at (0, -0.8) {\textbf{Bảng\_LuaChon}\\Nut, ChiSo\\Nhan, NutKeTiep};
  \node[tbl=purple] (smet) at (0, -2.6) {\textbf{Bảng\_SieuDuLieu}\\Chuong\_ID\\MaBam, NgayCapNhat};

  \node[grp=orange, minimum width=3.4cm, minimum height=4.0cm] (g3box) at (4.5, 2.0) {};
  \node[font=\scriptsize\bfseries\color{orange}, align=center] at (4.5, -4.85) {Nhóm 3 — Đám mây};
  \node[tbl=orange] (sync) at (4.5, 2.5) {\textbf{Bảng\_XepHang}\\TenNguoiDung, Diem\\ThoiGian, TocDo\\Truyen, Chuong};

  \node[grp=orange, minimum width=3.4cm, minimum height=1.2cm] (idbfs) at (4.5, -2.6) {\textbf{Hệ thống tệp ảo}\\(Chỉ trên Web)\\Đồng bộ sau mỗi lần ghi};

  \draw[arrow] (rec.north) -- ++(0,1.8) coordinate (top) -| (sync.north);
  \path (top) -- (top -| sync.north) node[midway, above, font=\tiny] {Khởi tạo Bảng xếp hạng};
  \draw[arrow](g2box.east |- g3box.west)--(g3box.west) node[midway,above,font=\tiny,align=center]{Đồng bộ D1};
  \draw[sdsArrow=orange,dashed](g1box.south)--++(0,-0.5)-|(idbfs.south);
\end{tikzpicture}
\end{adjustbox}
\end{center}

% ── CÔNG TÁC VẬN HÀNH (DEVOPS) ──────────────────────────────────────────────
\subsection{Công tác Vận hành \& Tự động hóa (DevOps)}

\subsubsection{Hệ thống Đóng gói Đa nền tảng}
Chỉ cần chạy một lệnh \texttt{build.sh} (macOS/Linux) hoặc \texttt{build.ps1} (Windows), hệ thống tự động kiểm tra, cài đặt thư viện thiếu và đóng gói cho 4 mục tiêu: Máy tính, Trình duyệt Web, và các Phân hệ độc lập để kiểm thử.

\begin{center}
\begin{adjustbox}{max width=\linewidth}
\begin{tikzpicture}[
  b/.style={techBoxSm=tOrange, text width=2.6cm},
  outBox/.style={techBoxSm=tGreen, text width=2.4cm},
  arr/.style={sdsArrow=tOrange}
]
  \node[techBox=tBlue, text width=6.5cm] (src) at (0,4.0) {Mã nguồn C++ \& Cấu hình CMake};
  
  \node[b](bs)  at (-3.5, 1.5) {Tập lệnh Tự động\\(Linux/Windows)};
  \node[b](cmake) at (0, 1.5) {Hệ thống\\Xây dựng (CMake)};
  \node[b](emsdk) at (3.5, 1.5) {Bộ công cụ Web\\(Emscripten)};

  \node[outBox](nat) at (-4.5, -1.5) {Tệp chạy\\Máy tính\\(.exe/.app)};
  \node[outBox](wasm) at (-1.5, -1.5) {Tệp Web\\(.html + .js + .wasm)};
  \node[outBox](sa)  at (1.5, -1.5)  {Các Tệp\\Kiểm thử\\Độc lập};
  \node[outBox](ci)  at (4.5, -1.5)  {Kết quả tự động\\(Máy chủ Tích hợp)};

  \draw[arr](src)--(bs); \draw[arr](src)--(cmake); \draw[arr](src)--(emsdk);
  \draw[arr](bs)--(nat); \draw[arr](cmake)--(nat); \draw[arr](cmake)--(sa);
  \draw[arr, rounded corners=12pt] (emsdk.east) -- ++(2.8,0) |- ([yshift=-0.8cm]wasm.south) -- (wasm.south);
  \draw[arr](cmake)--(wasm); \draw[arr](cmake)--(ci);
\end{tikzpicture}
\end{adjustbox}
\end{center}

\subsubsection{Sự hỗ trợ Đa nền tảng}
\begin{center}
\begin{adjustbox}{max width=\linewidth}
\begin{tabular}{|p{5.5cm}|c|c|c|c|}
  \hline
  \rowcolor{blue!15}
  \textbf{Chức năng} & \textbf{macOS} & \textbf{Windows} & \textbf{Linux} & \textbf{Web} \\
  \hline
  Vẽ Đồ họa Cửa sổ & Hỗ trợ & Hỗ trợ & Hỗ trợ & Hỗ trợ \\
  \hline
  Phát Luồng Âm thanh & Hỗ trợ & Hỗ trợ & Hỗ trợ & Một phần \\
  \hline
  Lưu trữ CSDL Cục bộ & Hỗ trợ & Hỗ trợ & Hỗ trợ & Một phần \\
  \hline
  Lưu trữ Tệp Ảo Web & Không & Không & Không & Hỗ trợ \\
  \hline
  Giao tiếp Mạng & Hỗ trợ & Hỗ trợ & Hỗ trợ & Hỗ trợ \\
  \hline
  Biểu tượng Gốc HĐH & Hỗ trợ & Hỗ trợ & Không & Không \\
  \hline
\end{tabular}
\end{adjustbox}
\end{center}

\subsubsection{Chuỗi Phụ thuộc (Dependencies)}
\begin{center}
\begin{adjustbox}{max width=\linewidth}
\begin{tikzpicture}[
  pkg/.style={techBoxSm=tOrange, text width=2.4cm},
  from/.style={techBoxSm=tBlue, text width=2.4cm},
  arr/.style={sdsArrow=gray!60},
]
  \node[from](sdl3){Thư viện SDL3\\(Tự biên dịch)};
  \node[pkg,right=0.5cm of sdl3](nano){Xử lý ảnh Vector\\(Tệp tiêu đề đơn)};
  \node[pkg,right=0.5cm of nano](sq){Cơ sở dữ liệu\\SQLite3};
  \node[pkg,right=0.5cm of sq](json){Phân tích JSON\\(Tệp tiêu đề đơn)};
  \node[pkg,right=0.5cm of json](curl){Mạng HTTP\\(Hệ thống / Web)};
  \node[techBox=tBlue, text width=13.5cm, below=0.8cm of sq]
    (app) {Tệp chạy Máy tính / Gói Tệp chạy Web};
    
  \foreach \x in {sdl3,nano,sq,json,curl} {\draw[arr](\x)--(app.north -| \x);}
\end{tikzpicture}
\end{adjustbox}
\end{center}

\subsubsection{Tích hợp Liên tục (CI/CD) - 4 Luồng Công Việc}
Hệ thống tự động hóa qua GitHub Actions bảo vệ và triển khai dự án qua 4 luồng tách biệt.

\begin{center}
\begin{adjustbox}{max width=\linewidth}
\begin{tabular}{|p{3.5cm}|p{3.5cm}|p{4.5cm}|p{3.5cm}|}
  \hline
  \rowcolor{orange!15}
  \textbf{Luồng công việc} & \textbf{Kích hoạt khi} & \textbf{Mục tiêu} & \textbf{Môi trường} \\
  \hline
  \textbf{Phát hành Web} & Đẩy mã (Phân hệ) & Biên dịch Web $\rightarrow$ Đưa lên Trang trực tuyến & Ubuntu/ Windows/ macOS \\
  \hline
  \textbf{Đồng bộ Truyện} & Đẩy mã (Truyện) & Xác thực $\rightarrow$ Đẩy lên Gist $\rightarrow$ Xuất Liên kết & Ubuntu \\
  \hline
  \textbf{Kiểm tra Lỗi} & Đẩy mã (Nhánh phụ) & Biên dịch Máy tính + Web trên cả 3 HĐH cùng lúc & Ma trận 3 HĐH \\
  \hline
  \textbf{Phát hành API} & Đẩy mã (Mạng) & Xác thực Khóa $\rightarrow$ Đưa lên Đám mây $\rightarrow$ Kiểm tra sức khỏe & Ubuntu \\
  \hline
\end{tabular}
\end{adjustbox}
\end{center}

\textbf{Luồng 1: Phát hành Trang Web (Web Deployment)}
\begin{center}
\begin{adjustbox}{max width=\linewidth}
\begin{tikzpicture}[
  job/.style={techBoxSm=tOrange, text width=2.6cm, minimum height=1.0cm},
  trig/.style={techBoxSm=tBlue, text width=2.4cm},
  outBox/.style={techBoxSm=tGreen, text width=2.4cm},
  cSide/.style={rectangle, draw=gray!50, fill=gray!6, rounded corners=2pt, font=\tiny, align=center, inner sep=3pt, text width=2.2cm},
  arr/.style={sdsArrow=tOrange}, darr/.style={sdsArrow=gray!50, dashed}
]
  \node[trig](trig) at (-5.5,0) {Đẩy mã lên\\Nhánh chính};
  \node[job](cfg) at (-2.6,0) {\textbf{Tác vụ 0}\\Xác định Môi trường};
  \node[job](bld) at (0.8,0) {\textbf{Tác vụ 1}\\Biên dịch Web \& Đẩy tệp lên};
  \node[job](dep) at (4.2,0) {\textbf{Tác vụ 2}\\Phát hành Trang};
  \node[outBox](url) at (4.2,-1.8) {Liên kết Trò chơi Trực tuyến};
  \node[cSide](cache) at (0.8,-1.8) {Bộ nhớ đệm:\\Tải Thư viện \& Công cụ};
  \node[cSide](art) at (0.8,1.6) {Báo cáo:\\Nhật ký biên dịch};

  \draw[arr](trig)--(cfg); \draw[arr](cfg)--(bld) node[above,font=\tiny,midway]{HĐH}; \draw[arr](bld)--(dep); \draw[arr](dep)--(url);
  \draw[darr](cache)--(bld); \draw[darr](bld)--(art);
\end{tikzpicture}
\end{adjustbox}
\end{center}

\textbf{Luồng 2: Đồng bộ Chương Truyện lên Đám Mây Gist}
\begin{center}
\begin{adjustbox}{max width=\linewidth}
\begin{tikzpicture}[
  cStep/.style={techBoxSm=tGreen, text width=2.8cm, minimum height=0.95cm},
  trig/.style={techBoxSm=tBlue, text width=2.4cm},
  cSide/.style={rectangle, draw=gray!50, fill=gray!6, rounded corners=2pt, font=\tiny, align=center, inner sep=3pt, text width=2.8cm},
  arr/.style={sdsArrow=tGreen}, darr/.style={sdsArrow=gray!50, dashed}
]
  \node[trig](trig) at (-5.2,0) {Đẩy mã Truyện\\hoặc Thủ công};
  \node[cStep](s0) at (-2.0,0) {\textbf{Bước 0}\\Xác thực Khóa API};
  \node[cStep](s1) at (1.2,0) {\textbf{Bước 1}\\Đồng bộ Chương};
  \node[cStep](s2) at (4.4,0) {\textbf{Bước 2}\\Tái tạo tệp Khai báo};
  \node[cStep](s3) at (7.6,0) {\textbf{Bước 3}\\In Liên kết tải truyện};
  \node[cSide](ver) at (7.6,-1.8) {Cập nhật Liên kết vào Cấu hình ứng dụng};
  \node[cSide](sha) at (-2.0,-1.8) {Chặn lỗi từ đầu nếu Khóa không đủ quyền};
  
  \draw[arr](trig)--(s0); \draw[arr](s0)--(s1); \draw[arr](s1)--(s2); \draw[arr](s2)--(s3);
  \draw[darr](s3)--(ver); \draw[darr](s0)--(sha);
\end{tikzpicture}
\end{adjustbox}
\end{center}

\textbf{Luồng 3: Kiểm tra Biên dịch Đa nền tảng}
\begin{center}
\begin{adjustbox}{max width=\linewidth}
\begin{tikzpicture}[
  cStep/.style={techBoxSm=tPurple, text width=2.4cm, minimum height=0.95cm},
  os/.style={rectangle, draw=tPurple, fill=tPurple!8, rounded corners=3pt, font=\scriptsize\bfseries, align=center, inner sep=5pt, minimum width=2.0cm, minimum height=0.7cm},
  trig/.style={techBoxSm=tBlue, text width=2.4cm},
  arr/.style={sdsArrow=tPurple}, darr/.style={sdsArrow=gray!50, dashed}
]
  \node[trig](trig) at (0,0) {Đẩy mã lên\\Nhánh Phụ};
  \node[os](win)  at (3.0, 1.2) {Windows};
  \node[os](ubu)  at (3.0, 0.0) {Ubuntu};
  \node[os](mac)  at (3.0,-1.2) {macOS};

  \node[cStep](setup) at (5.8,0) {\textbf{Bước 1}\\Cài Thư viện};
  \node[cStep](nat)   at (8.6,0) {\textbf{Bước 2}\\Biên dịch Máy tính};
  \node[cStep](wasm)  at (11.4,0) {\textbf{Bước 3}\\Biên dịch Web};
  \node[rectangle, draw=gray!50, fill=gray!6, rounded corners=2pt, font=\tiny, align=center, inner sep=3pt, text width=2.2cm, right=0.5cm of wasm] (res) {Báo cáo \&\\Lưu Nhật ký lỗi};

  \foreach \os in {win,ubu,mac}{ \draw[arr](trig)--(\os); \draw[arr](\os)--(setup); }
  \draw[arr](setup)--(nat); \draw[arr](nat)--(wasm); \draw[darr](wasm)--(res);
\end{tikzpicture}
\end{adjustbox}
\end{center}

\textbf{Luồng 4: Triển khai Máy chủ Đám mây Cloudflare}
\begin{center}
\begin{adjustbox}{max width=\linewidth}
\begin{tikzpicture}[
  cStep/.style={techBoxSm=tOrange, text width=2.4cm, minimum height=0.95cm},
  trig/.style={techBoxSm=tBlue, text width=2.4cm},
  cSide/.style={rectangle, draw=gray!50, fill=gray!6, rounded corners=2pt, font=\tiny, align=center, inner sep=3pt, text width=2.4cm},
  arr/.style={sdsArrow=tOrange}, darr/.style={sdsArrow=gray!50, dashed}
]
  \node[trig](trig) at (0,0) {Đẩy mã Mạng\\(API)};
  \node[cStep](s0) at (3.2,0) {\textbf{Bước 0}\\Xác thực Khóa};
  \node[cStep](s1) at (6.4,0) {\textbf{Bước 1}\\Đẩy mã Máy chủ};
  \node[cStep](s2) at (9.6,0) {\textbf{Bước 2}\\Kiểm tra Sức khỏe};
  \node[cSide](info) at (9.6,-1.8) {Gọi API xem hoạt động chưa, xuất thông báo};
  
  \draw[arr](trig)--(s0); \draw[arr](s0)--(s1); \draw[arr](s1)--(s2); \draw[darr](s2)--(info);
\end{tikzpicture}
\end{adjustbox}
\end{center}

\label{LastPageTechnologyAdoption}
\end{document}